\documentclass[11pt,a4paper]{article}

\usepackage[margin=1in]{geometry}
\usepackage{graphicx}
\usepackage{subcaption}
\usepackage{booktabs}
\usepackage{hyperref}
\usepackage{xcolor}

\hypersetup{
  colorlinks=true,
  linkcolor=blue,
  citecolor=blue,
  urlcolor=blue
}

\newcommand{\cms}{CMS}
\newcommand{\figwidth}{0.46\textwidth}
\newcommand{\plotplaceholder}[2]{%
  \fbox{\begin{minipage}[c][0.26\textheight][c]{#1}
  \centering
  \small Missing final production figure:\\
  \scriptsize\nolinkurl{#2}
  \end{minipage}}%
}
\newcommand{\maybeplot}[2]{%
  \IfFileExists{#2}{\includegraphics[width=#1]{#2}}{\plotplaceholder{#1}{#2}}%
}

\title{\bfseries NanoAOD Lepton Reweighting and Validation}
\author{CMS Simulation Study}
\date{\today}

\begin{document}
\maketitle

\begin{abstract}
This note documents a small NanoAOD-level workflow used to distort lepton
kinematics and identification efficiencies in simulated events.  The modified
NanoAOD files are compared to the raw inputs through one-dimensional
distributions, efficiency curves, tag-and-probe mass fits, and two-dimensional
correlation maps.  The final figures in this note must be produced from the
full set of input ROOT files after rerunning \texttt{modify\_nanoaod} and
\texttt{plot.py}.
\end{abstract}

\section{Workflow}

The C++ program \texttt{modify\_nanoaod.cpp} reads NanoAOD \texttt{Events}
trees, applies deterministic lepton momentum scale and resolution changes, and
distorts selected lepton ID/isolation working points according to the
configuration in \texttt{config.json}.  The modified files are written to the
\texttt{modified/} directory with the configured suffix.

The Python program \texttt{plot.py} compares each raw input file to its
corresponding modified output.  It produces normalized lepton distributions,
dilepton spectra, MC-truth and tag-and-probe efficiency curves, fit diagnostic
plots, and two-dimensional heatmaps.  For every input/output heatmap pair, an
additional \texttt{Modified - raw} map is produced.

\section{Configuration}

The lepton momentum scale, resolution, and efficiency distortions use moderate
configured amplitudes.  The barrel and endcap scale shifts are kept in the
same direction with small regional differences, and the efficiency turn-on
distortions are limited, so the validation changes remain visible without
splitting the dilepton invariant-mass peak into artificial multiple peaks.
Branches with explicit energy values are scaled by the same ratio as the
corresponding transverse momentum, so the energy response follows the
configured momentum distortion without introducing an independent mass-shape
change.

For tag-and-probe fits, the nominal invariant-mass fit range is kept broad
($55$--$135$ GeV).  For bins with an additional high-mass background shape,
the fit window can be shortened to 100 GeV, while the signal coefficients are
kept normalized over the full mass range.  The reported efficiency therefore
corresponds to the full signal PDF normalization, not only to the restricted
fit window.  Each accepted fit must satisfy RooFit status and covariance
quality requirements and pass a pass/fail $\chi^{2}/\mathrm{ndf}$ validation.
The $\chi^{2}/\mathrm{ndf}$ values and the fit window are printed directly on
each fit PDF.

\section{Kinematic Distributions}

Figure~\ref{fig:lepton-kinematics} shows representative lepton-level
distributions before and after the NanoAOD modification.  The relative momentum
shift distributions directly expose the configured scale changes, while the
smearing response isolates the stochastic resolution term.

\begin{figure}[htbp]
  \centering
  \begin{subfigure}{\figwidth}
    \centering
    \maybeplot{\linewidth}{../figures/distributions/lepton/dist_muon_pt.pdf}
    \caption{Muon $p_{\mathrm{T}}$}
  \end{subfigure}
  \begin{subfigure}{\figwidth}
    \centering
    \maybeplot{\linewidth}{../figures/distributions/lepton/dist_electron_pt.pdf}
    \caption{Electron $p_{\mathrm{T}}$}
  \end{subfigure}
  \begin{subfigure}{\figwidth}
    \centering
    \maybeplot{\linewidth}{../figures/distributions/lepton/dist_muon_relative_pt_shift.pdf}
    \caption{Muon relative $p_{\mathrm{T}}$ shift}
  \end{subfigure}
  \begin{subfigure}{\figwidth}
    \centering
    \maybeplot{\linewidth}{../figures/distributions/lepton/dist_electron_relative_pt_shift.pdf}
    \caption{Electron relative $p_{\mathrm{T}}$ shift}
  \end{subfigure}
  \caption{Representative lepton kinematic validation plots.}
  \label{fig:lepton-kinematics}
\end{figure}

\section{Efficiency Validation}

The efficiency validation compares the raw and modified samples with both
generator-matched MC truth and tag-and-probe fits.  The modified efficiencies
are expected to follow the configured global and kinematic distortions.  Fit
diagnostics are kept under \texttt{figures/efficiency/fit/}; rejected fits are
stored separately and their efficiency points fall back to counting estimates.

\begin{figure}[htbp]
  \centering
  \begin{subfigure}{\figwidth}
    \centering
    \maybeplot{\linewidth}{../figures/efficiency/eff_muon_Muon_mediumId_pt.pdf}
    \caption{Muon medium ID vs. $p_{\mathrm{T}}$}
  \end{subfigure}
  \begin{subfigure}{\figwidth}
    \centering
    \maybeplot{\linewidth}{../figures/efficiency/eff_electron_Electron_mvaFall17V2Iso_WP80_pt.pdf}
    \caption{Electron WP80 vs. $p_{\mathrm{T}}$}
  \end{subfigure}
  \caption{Representative efficiency curves.}
  \label{fig:efficiency}
\end{figure}

\section{Dilepton Resolution}

The dilepton mass distributions summarize the combined impact of the lepton
momentum scale and resolution changes.  The modified mass peak can shift or
broaden relative to the raw sample according to the configured scale and
resolution terms.  The validation should therefore check that the change is
controlled and does not introduce artificial multi-peak mass shapes.

\begin{figure}[htbp]
  \centering
  \begin{subfigure}{\figwidth}
    \centering
    \maybeplot{\linewidth}{../figures/distributions/dilepton/dilepton_muon_mass.pdf}
    \caption{Muon-pair mass}
  \end{subfigure}
  \begin{subfigure}{\figwidth}
    \centering
    \maybeplot{\linewidth}{../figures/distributions/dilepton/dilepton_electron_mass.pdf}
    \caption{Electron-pair mass}
  \end{subfigure}
  \caption{Dilepton mass validation.}
  \label{fig:dilepton}
\end{figure}

\section{Two-dimensional Heatmaps}

The correlation plots are written in separate \texttt{input}, \texttt{output},
and \texttt{delta} directories.  The delta maps show the bin-by-bin
\texttt{Modified - raw} entry difference and provide a compact check of where
the modification changes the event population.

\begin{figure}[htbp]
  \centering
  \begin{subfigure}{\figwidth}
    \centering
    \maybeplot{\linewidth}{../figures/correlation/input/corr_muon_input_energy_pt.pdf}
    \caption{Muon raw energy vs. $p_{\mathrm{T}}$}
  \end{subfigure}
  \begin{subfigure}{\figwidth}
    \centering
    \maybeplot{\linewidth}{../figures/correlation/delta/corr_muon_delta_energy_pt.pdf}
    \caption{Muon modified minus raw}
  \end{subfigure}
  \begin{subfigure}{\figwidth}
    \centering
    \maybeplot{\linewidth}{../figures/correlation/input/corr_electron_input_energy_pt.pdf}
    \caption{Electron raw energy vs. $p_{\mathrm{T}}$}
  \end{subfigure}
  \begin{subfigure}{\figwidth}
    \centering
    \maybeplot{\linewidth}{../figures/correlation/delta/corr_electron_delta_energy_pt.pdf}
    \caption{Electron modified minus raw}
  \end{subfigure}
  \caption{Representative two-dimensional heatmaps.}
  \label{fig:heatmaps}
\end{figure}

\section{Production Status}

The final production requires the CMS and LCG CVMFS environments for
\texttt{modify\_nanoaod}.  The command sequence is captured in
\texttt{scripts/run\_modify\_all\_cvmfs.sh}.  Plotting is performed in the pixi
environment with \texttt{scripts/run\_plot\_all\_pixi.sh}.  This note should be
rebuilt after both production steps complete successfully.

\end{document}
